% !PRINTED 72
\weberpage{72}
%<<CONT

sind, nichts ändern, so dass für \(\Delta\) nur der Werth Null übrig bleibt.

Man drückt den Satz V. nur anders aus, wenn man sagt:

\gesperrt{VI. Man erhält eine verschwindende Determinante, wenn man die Elemente einer Reihe durch die entsprechenden Elemente einer anderen Reihe, oder, kurz gesagt, wenn man einen unteren oder oberen Index durch einen anderen ersetzt.}

\begin{center}
{\large \S.\ 22.}

\vspace{0.5em}
{\large\gesperrt{Unterdeterminanten.}}
\end{center}

In jedem Gliede der entwickelten Determinante
\[
\sum \pm a_1^{(1)}a_2^{(2)}\cdots a_n^{(n)},
\]
deren Werth wir jetzt mit \(A\) bezeichnen wollen, kommt jede der Zahlen \(1,2,\ldots,n\) ein und nur einmal als unterer Index vor. Es wird also ein gewisser Complex von Gliedern den Factor \(a_1^{(1)}\) enthalten, ein anderer Complex den Factor \(a_1^{(2)}\) u.~s.~f., endlich ein Complex den Factor \(a_1^{(n)}\); jedes Glied der Determinante kommt in einem und nur in einem dieser Complexe vor.

Bezeichnen wir also den ersten dieser Complexe mit \(a_1^{(1)}A_1^{(1)}\), den zweiten mit \(a_1^{(2)}A_1^{(2)}\), den letzten mit \(a_1^{(n)}A_1^{(n)}\), so können wir die Determinante folgendermaassen darstellen:
\begin{equation}
A=a_1^{(1)}A_1^{(1)}+a_1^{(2)}A_1^{(2)}+\cdots+a_1^{(n)}A_1^{(n)}.
\tag{1}
\end{equation}
An Stelle des unteren Index \(1\) hätten wir ebenso gut jeden anderen, \(\nu\), herausgreifen und daher
\begin{equation}
A=a_\nu^{(1)}A_\nu^{(1)}+a_\nu^{(2)}A_\nu^{(2)}+\cdots+a_\nu^{(n)}A_\nu^{(n)}
\tag{2}
\end{equation}
setzen können. Darin bedeutet das Product \(a_\nu^{(\mu)}A_\nu^{(\mu)}\) den Complex aller Glieder der Determinante, die den Factor \(a_\nu^{(\mu)}\) enthalten.

Da dieselben Regeln wie für die unteren so auch für die oberen Indices gelten, so kann man die Determinante auch noch in der folgenden Weise schreiben:
\begin{equation}
A=a_1^{(\mu)}A_1^{(\mu)}+a_2^{(\mu)}A_2^{(\mu)}+\cdots+a_n^{(\mu)}A_n^{(\mu)},
\tag{3}
\end{equation}
worin \(\mu\) gleichfalls jeden der Indices \(1,2,\ldots,n\) bedeuten kann.
